\documentclass[11pt]{article}
\usepackage{mathblatt}
\usepackage{multicol}


\begin{document}
\pfheftstil
\pfheftkopf{Quadratische \textperiodcentered{} Lösungen}{}
\begin{pfloesung}{}
\lz{1.}{x $-$ ($-$3), d = $-$3; x $-$ 2, d = 2; x $-$ ($-$1), d = $-$1}{}{}
\end{pfloesung}
\begin{pfloesung}{}
\lz{2.}{x² $-$ 6x + 9; x² + 6x + 9; x² $-$ 4x + 4}{}{}
\end{pfloesung}
\begin{pfloesung}{}
\lz{3.}{$-$6 | 7 | 2; 4 | 3 | 1; $-$2 | $-$3 | 4}{}{}
\end{pfloesung}
\begin{pfloesung}{}
\lz{4.}{x² + x $-$ 6 = 0; x² + 2x $-$ 3 = 0}{}{}
\end{pfloesung}
\begin{pfloesung}{}
\lz{5.}{$4$}{Scheitel bei x}{}
\end{pfloesung}
\begin{pfloesung}{}
\lz{6.}{Normalform}{}{}
\end{pfloesung}
\begin{pfloesung}{}
\lz{7.}{$S(2|5)$}{}{}
\end{pfloesung}
\begin{pfloesung}{}
\lz{8.}{S($-$3 | $-$2)}{}{}
\end{pfloesung}
\begin{pfloesung}{}
\lz{9.}{S(2 | $-$2)}{}{}
\end{pfloesung}
\begin{pfloesung}{}
\lz{10.}{S($-$2 | $-$4)}{}{}
\end{pfloesung}
\begin{pfloesung}{}
\lz{11.}{S(2|$-$4)}{2, e = $-$4}{}
\end{pfloesung}
\begin{pfloesung}{}
\lz{12.}{S($-$2 | $-$1)}{}{}
\end{pfloesung}
\begin{pfloesung}{}
\lz{13.}{$(x - 5)^2 + 2$}{f(x)}{}
\end{pfloesung}
\begin{pfloesung}{}
\lz{14.}{y = (x $-$ 1)² + 3}{}{}
\end{pfloesung}
\begin{pfloesung}{}
\lz{15.}{y = 3x² $-$ 2}{}{}
\end{pfloesung}
\begin{pfloesung}{}
\lz{16.}{$(x + 4)^2 + 2$, Scheitel $S(-4|2)$}{g(x)}{}
\end{pfloesung}
\begin{pfloesung}{}
\lz{17.}{S(2 | 0)}{}{}
\end{pfloesung}
\begin{pfloesung}{}
\lz{18.}{p(x) = (x $-$ 2)² $-$ 2}{x² $-$ 4x + 4 $-$ 4 + 2}{}
\end{pfloesung}
\begin{pfloesung}{}
\lz{19.}{Skizze: Normalparabel mit Scheitel (2 | $-$4); y = (x $-$ 2)² $-$ 4}{Scheitel: ($-$2 | $-$4) = (2 | $-$4)}{}
\end{pfloesung}
\begin{pfloesung}{}
\lz{20.}{y = (x + 1)² $-$ 5 $\Rightarrow$ S($-$1; $-$5)}{}{}
\end{pfloesung}
\begin{pfloesung}{}
\lz{21.}{y = $-$(x $-$ 2)² + 3}{}{}
\end{pfloesung}
\begin{pfloesung}{}
\lz{22.}{y = x² + 2}{}{}
\end{pfloesung}
\begin{pfloesung}{}
\lz{23.}{q: y = $-$(x + 3)²}{}{}
\end{pfloesung}
\begin{pfloesung}{}
\lz{24.}{S(3|$-$2); p(x) = (x $-$ 3)² $-$ 2}{x² $-$ 6x + 9 $-$ 2 (quadratische Ergänzung)}{}
\end{pfloesung}
\begin{pfloesung}{}
\lz{25.}{S($-$1|6); p(x) = $-$(x + 1)² + 6}{$-$1, e = 6}{}
\end{pfloesung}
\begin{pfloesung}{}
\lz{26.}{Skizze: Normalparabel mit Scheitel S$_y$(0 | e), e > 0; z. B. q(x) = x² + 1; z. B. q$'$(x) = $-$x² $-$ 1}{e > 0, nach oben geöffnet}{}
\end{pfloesung}
\begin{pfloesung}{}
\lz{27.}{Formel}{Normalform}{}
\end{pfloesung}
\begin{pfloesung}{}
\lz{28.}{$-12$}{Prozentsatz: -1, q}{}
\end{pfloesung}
\begin{pfloesung}{}
\lz{29.}{$x + 4$}{$x^2 - 2x$}{}
\end{pfloesung}
\begin{pfloesung}{}
\lz{30.}{x$_1$ = $-$1, x$_2$ = $-$21}{}{}
\end{pfloesung}
\begin{pfloesung}{}
\lz{31.}{x$_1$ = $-$10; x$_2$ = 10}{}{}
\end{pfloesung}
\begin{pfloesung}{}
\lz{32.}{nein}{0,5(x $-$ 3)² = 0 bzw. Diskriminante 9 $-$ 9 = 0, nur x = 3}{}
\end{pfloesung}
\begin{pfloesung}{}
\lz{33.}{6x = x² + 9 $\Rightarrow$ (x $-$ 3)² = 0 $\Rightarrow$ x = 3}{}{}
\end{pfloesung}
\begin{pfloesung}{}
\lz{34.}{x$_1$ = $-$1; x$_2$ = $-$3}{2x² + 8x + 6 $\Rightarrow$ 0 = x² + 4x + 3 = 0; $-$2 ± 1}{}
\end{pfloesung}
\begin{pfloesung}{}
\lz{35.}{x² $-$ 80x + 8 786 = 0; Radikand 1 600 $-$ 8 786 = $-$7 186 < 0 – keine Nullstellen}{}{}
\end{pfloesung}
\begin{pfloesung}{}
\lz{36.}{Höhe 30 cm; x² = 600 $\Rightarrow$ x $\approx$ ±24,5; Weite $\approx$ 49 cm}{}{}
\end{pfloesung}
\begin{pfloesung}{}
\lz{37.}{0 = x² $-$ 4x + 1 $\Rightarrow$ x$_1$ $\approx$ 0,27; x$_2$ $\approx$ 3,73; Breite $\approx$ 3,46 m}{}{}
\end{pfloesung}
\begin{pfloesung}{}
\lz{38.}{N1N$_1$(1,59|0), N2N$_2$(4,41|0), x = 3 ± √2}{Diskriminante: 9 $-$ 7 $\Rightarrow$ 2; √2 $\approx$ 1,41}{}
\end{pfloesung}
\begin{pfloesung}{}
\lz{39.}{(x + 3)² $-$ 2 = x² + 6x + 9 $-$ 2 = x² + 6x + 7; x$_1$ = $-$3 + √2 $\approx$ $-$1,59, x$_2$ = $-$3 $-$ √2 $\approx$ $-$4,41}{(x + 3)² $\Rightarrow$ 2 = $-$3 ± √2}{}
\end{pfloesung}
\begin{pfloesung}{}
\lz{40.}{Q(3; $-$3)}{}{}
\end{pfloesung}
\begin{pfloesung}{}
\lz{41.}{S$_1$($-$1 | $-$2); S$_2$(2 | 7)}{x² + 2x $-$ 1 $\Rightarrow$ 3x + 1 = x² $-$ x $-$ 2 = 0; $-$1, x$_2$ = 2; 3x + 1 $\Rightarrow$ $-$2, y$_2$ = 7}{}
\end{pfloesung}
\begin{pfloesung}{}
\lz{42.}{S$_1$($-$1|5), S$_2$(3|$-$3)}{x² $-$ 2x $-$ 3 $\Rightarrow$ 0; 1 ± 2}{}
\end{pfloesung}
\begin{pfloesung}{}
\lz{43.}{S1($-$1|$-$3), S2(5|21)}{x² $-$ 4x $-$ 5 $\Rightarrow$ 0; 2 ± 3}{}
\end{pfloesung}
\begin{pfloesung}{}
\lz{44.}{Gerade}{}{}
\end{pfloesung}
\begin{pfloesung}{}
\lz{45.}{erste Tabelle (y = 12, 3, 0, 3, 12)}{3 · 2² $\Rightarrow$ 12}{}
\end{pfloesung}
\begin{pfloesung}{}
\lz{46.}{h(t) = $-$3t² + 2400; h(0) = 2400 und Parabel nach unten geöffnet $\Rightarrow$ $-$3t² + 2400}{h(20) = $-$3 · 20² + 2400 = 1200}{}
\end{pfloesung}
\begin{pfloesung}{}
\lz{47.}{$36$}{f(-6) = $(-6)^2$}{}
\end{pfloesung}
\begin{pfloesung}{}
\lz{48.}{Nein}{p($-$5) = 25 $-$ 4 = 21 $\neq$ 20; ($-$5)² = 25}{}
\end{pfloesung}
\begin{pfloesung}{}
\lz{49.}{D($-$1,5 | $-$2,5)}{}{}
\end{pfloesung}
\begin{pfloesung}{}
\lz{50.}{y = 12}{(2 + 2)² $-$ 4}{}
\end{pfloesung}
\begin{pfloesung}{}
\lz{51.}{wahr; falsch}{p($-$1) = 7; p(0) = 2 = Schnittpunkt (0 | 2)}{}
\end{pfloesung}
\begin{pfloesung}{}
\lz{52.}{ja}{S($-$3 | $-$9) ist Schnittpunkt; p($-$3) = $-$($-$3)² = $-$9; g($-$3) = 2 · ($-$3) $-$ 3 = $-$9}{}
\end{pfloesung}
\begin{pfloesung}{}
\lz{53.}{Normalparabel mit Scheitel ($-$1; $-$4), Nullstellen $-$3 und 1}{}{}
\end{pfloesung}
\begin{pfloesung}{}
\lz{54.}{Bogen durch $(-2|16)$, $(-1|4)$, $(0|0)$, $(1|4)$, $(2|16)$}{}{}
\end{pfloesung}
\begin{pfloesung}{}
\lz{55.}{Parabel mit Scheitel (0|$-$4) durch (±2|0) und (±3|5); S(0|$-$4)}{Normalparabel um: 4 nach unten verschoben; Nullstellen ±2}{}
\end{pfloesung}
\begin{pfloesung}{}
\lz{56.}{$c$ mit $c$ kleiner als $-2$, denn der Scheitel $S(0|{-2})$ ist der tiefste Punkt}{-3; jede waagerechte Gerade g(x)}{}
\end{pfloesung}
\begin{pfloesung}{}
\lz{57.}{kein gemeinsamer Punkt: Scheitel $(1|2)$ und $(1|{-4})$ übereinander, gleiche Öffnung}{g ist f um: 6 nach unten verschoben}{}
\end{pfloesung}
\begin{pfloesung}{}
\lz{58.}{z. B. f(x) = 1}{Scheitel S(0 | 0), Parabel nach unten geöffnet}{}
\end{pfloesung}
\begin{pfloesung}{}
\lz{59.}{x kommt im Quadrat vor}{quadratisch}{}
\end{pfloesung}
\begin{pfloesung}{}
\lz{60.}{rechts positiv}{zwei Lösungen}{}
\end{pfloesung}
\begin{pfloesung}{}
\lz{61.}{die Vorzahl von $x^2$}{durch: 4}{}
\end{pfloesung}
\begin{pfloesung}{}
\lz{62.}{rechts steht null}{gilt}{}
\end{pfloesung}
\begin{pfloesung}{}
\lz{63.}{$x_1 = 9$, $x_2 = -9$}{}{}
\end{pfloesung}
\begin{pfloesung}{}
\lz{64.}{$25$ : ja}{$(-7 + 2)^2$ $\Rightarrow$ $(-5)^2$; wA}{}
\end{pfloesung}
\begin{pfloesung}{}
\lz{65.}{$-14$}{$-7,$ denn $(-7) \cdot (-7 + 9)$ $\Rightarrow$ $(-7) \cdot 2$}{}
\end{pfloesung}
\begin{pfloesung}{}
\lz{66.}{$0$ : ja}{$(-2)^2 - 3 \cdot (-2) - 10$ $\Rightarrow$ 4 + 6 - 10; wA}{}
\end{pfloesung}
\begin{pfloesung}{}
\lz{67.}{x = $-$2}{($-$2) · 3 $\Rightarrow$ $-$6}{}
\end{pfloesung}
\begin{pfloesung}{}
\lz{68.}{$x_1 = 2$, $x_2 = -3$}{$x^2 + x + 1$ $\Rightarrow$ $7, x^2 + x - 6$ = 0; $-0{,}5 \pm \sqrt{0{,}25 + 6}$ $\Rightarrow$ $-0{,}5 \pm 2{,}5$}{}
\end{pfloesung}
\begin{pfloesung}{}
\lz{69.}{$x_1 = 5$, $x_2 = -4$}{$2x^2 - 2x - 40$ $\Rightarrow$$0 \mid : 2$; $x^2 - x - 20$ $\Rightarrow$ 0; $x_{1,2}$ = $0{,}5 \pm \sqrt{0{,}25 + 20}$ = $0{,}5 \pm 4{,}5$}{}
\end{pfloesung}
\begin{pfloesung}{}
\lz{70.}{x$_1$ = $-$5, x$_2$ = 3}{x² + 2x $-$ 15 $\Rightarrow$ 0; $-$1 ± 4}{}
\end{pfloesung}
\begin{pfloesung}{}
\lz{71.}{1~Nullstellen berechnen, 2~x zu gegebenem y, 3~Nullstellen berechnen, 4~x zu gegebenem y, 5~Nullstellen berechnen}{}{}
\end{pfloesung}
\begin{pfloesung}{\textbf{Prüfstein}}
\lz{a)}{Gerade durch (0|2) und (1|0), fallend}{Steigungsdreieck: 1 nach rechts, 2 nach unten}{}
\lz{b)}{Ja}{f($-$4) = $-$2 · ($-$4) + 2 = 10; f($-$4) = 10}{}
\lz{c)}{S(2|$-$2)}{2, e = $-$2}{}
\lz{d)}{genau ein gemeinsamer Punkt S(2|$-$2), gleicher Scheitel, p nach oben, q nach unten geöffnet $\Rightarrow$ Berührung nur im Scheitel}{Scheitel beider Parabeln: (2|$-$2)}{}
\end{pfloesung}
\end{document}